\documentclass[12pt,a4paper]{article}
\usepackage[utf8]{inputenc}
\usepackage[T1]{fontenc}
\usepackage{babel}
\usepackage{hyperref}
\usepackage{listings}
\usepackage{xcolor}
\usepackage{array}
\usepackage[margin=2.5cm]{geometry}

\definecolor{codegreen}{rgb}{0,0.6,0}
\definecolor{codegray}{rgb}{0.5,0.5,0.5}
\definecolor{codepurple}{rgb}{0.58,0,0.82}
\definecolor{backcolour}{rgb}{0.95,0.95,0.92}
\definecolor{keyword}{rgb}{0,0,1}
\definecolor{string}{rgb}{0.58,0,0.82}
\definecolor{comment}{rgb}{0.57,0.57,0.57}
\definecolor{type}{rgb}{0.1,0.5,0.1}

\lstdefinestyle{cstyle}{
    backgroundcolor=\color{backcolour},
    commentstyle=\color{comment}\textit,
    keywordstyle=\color{keyword}\bfseries,
    stringstyle=\color{string},
    numberstyle=\tiny\color{codegray},
    basicstyle=\ttfamily\small,
    breaklines=true,
    frame=single,
    framesep=5pt,
    frameround=tttt,
    numbers=left,
    numbersep=8pt,
    captionpos=b,
    tabsize=4,
    showstringspaces=false
}

\lstset{style=cstyle}

\title{Implémentation des Widgets GTK4 \\ Rapport Technique - Rapport 2}
\author{Projet GTK Wrapper}
\date{\today}

\begin{document}

\maketitle

\tableofcontents

\newpage

\section{Introduction}
Ce rapport documente l'implémentation interne du wrapper GTK4. Contrairement au premier rapport, centré sur les structures de configuration et l'usage de l'API, ce second document décrit la logique des fonctions, les macros et les appels GTK réellement utilisés et les choix techniques effectués dans chaque module.

Les extraits de code présentés ici reprennent l'implémentation réelle du projet. Les commentaires ont été retirés dans les listings afin de conserver une compilation LaTeX simple et robuste tout en gardant un style visuel identique au premier rapport.

\section{Philosophie de conception : Approche Déclarative et C99}

L'une des innovations majeures de ce wrapper réside dans l'adoption d'un paradigme déclaratif au sein d'un langage procédural comme le C. Au lieu de contraindre le développeur à écrire de multiples appels impératifs (tels que \texttt{gtk\_window\_set\_title()}, \texttt{gtk\_window\_set\_default\_size()}, etc.), l'API centralise l'état désiré dans des structures de configuration (ex: \texttt{window\_config}, \texttt{button\_config}).

Cette approche est rendue particulièrement puissante et expressive grâce à l'utilisation des \textbf{initialisateurs désignés (Designated Initializers)} introduits par la norme C99. Cette fonctionnalité permet d'initialiser une structure en nommant explicitement les champs concernés :

\begin{lstlisting}[language=C]
window_config config = {
    .title = "Mon Application",
    .default_width = 800,
    .resizable = true
};
\end{lstlisting}

L'utilisation des initialisateurs désignés apporte trois bénéfices fondamentaux à l'architecture du projet :
\begin{enumerate}
    \item \textbf{Lisibilité accrue (Simulation d'arguments nommés) :} Contrairement aux constructeurs C traditionnels nécessitant de passer de nombreux arguments positionnels dont le sens se perd à la lecture, cette syntaxe documente le code de manière inhérente. Le développeur comprend immédiatement quelle propriété est modifiée, s'approchant de l'ergonomie des langages modernes (comme les paramètres nommés en Python ou les objets de configuration en JavaScript).
    \item \textbf{Sécurité par défaut (Zero-initialization) :} La norme C99 garantit que tout champ non mentionné explicitement dans un bloc d'initialisation désigné est automatiquement initialisé à zéro (les pointeurs deviennent \texttt{NULL}, les entiers \texttt{0}, et les booléens \texttt{false}). Le wrapper exploite cette garantie pour appliquer des valeurs par défaut sûres. Si un développeur omet le champ \texttt{.title}, le wrapper détectera un pointeur \texttt{NULL} et laissera GTK gérer le comportement par défaut, évitant ainsi les plantages liés aux valeurs résiduelles en mémoire.
    \item \textbf{Flexibilité et concision :} Les développeurs ne renseignent que les propriétés pertinentes pour leur contexte métier, sans avoir à fournir des valeurs nulles ou factices pour les paramètres dont ils n'ont pas l'utilité.
\end{enumerate}

\newpage

\section{Module \texttt{window.c}}

Le module \texttt{window.c} construit les fenêtres principales de l'application. Il encapsule la création d'un \texttt{GtkApplicationWindow}, l'application des propriétés usuelles et la gestion d'un fond personnalisé par CSS.

\subsection{Fonction \texttt{apply\_window\_background}}

Cette fonction privée applique un style CSS ciblé sur une seule fenêtre. L'idée est de nommer le widget puis de générer dynamiquement une règle CSS qui peut définir soit une couleur de fond seule, soit une couleur combinée à une image.

\begin{lstlisting}[language=C]
static void apply_window_background(const window_config *config, GtkWidget *window) {
    if (config == NULL || window == NULL || config->background_color == NULL || config->background_color[0] == '\0') {
        return;
    }

    const char *name = config->widget_name;
    if (name == NULL || name[0] == '\0') {
        name = "wrapper-window";
        gtk_widget_set_name(window, name);
    }

    GtkCssProvider *provider = gtk_css_provider_new();
    char *css;

    if (config->background_image_path != NULL && config->background_image_path[0] != '\0') {
        css = g_strdup_printf(
            "window#%s, window#%s > * {"
            "background-color: %s;"
            "background-image: url('%s');"
            "background-size: cover;"
            "background-position: center;"
            "}",
            name,
            name,
            config->background_color,
            config->background_image_path
        );
    } else {
        css = g_strdup_printf(
            "window#%s, window#%s > * { background-color: %s; }",
            name,
            name,
            config->background_color
        );
    }

    gtk_css_provider_load_from_string(provider, css);
    gtk_style_context_add_provider_for_display(
        gtk_widget_get_display(window),
        GTK_STYLE_PROVIDER(provider),
        GTK_STYLE_PROVIDER_PRIORITY_APPLICATION
    );

    g_free(css);
    g_object_unref(provider);
}
\end{lstlisting}

\subsection{Fonction \texttt{create\_window}}

La fonction principale du module instancie la fenêtre, applique les propriétés déclarées dans \texttt{window\_config}, délègue la personnalisation visuelle à \texttt{apply\_window\_background} puis applique le style commun partagé avec les autres widgets.

\begin{lstlisting}[language=C]
GtkWidget *create_window(GtkApplication *app, const window_config *config) {
    GtkWidget *window = gtk_application_window_new(app);
    if (config == NULL) {
        return window;
    }

    if (config->title) {
        gtk_window_set_title(GTK_WINDOW(window), config->title);
    }
    if (config->icon_name) {
        gtk_window_set_icon_name(GTK_WINDOW(window), config->icon_name);
    }
    if (config->widget_name) {
        gtk_widget_set_name(window, config->widget_name);
    }

    gtk_window_set_default_size(GTK_WINDOW(window), config->default_width, config->default_height);
    gtk_window_set_resizable(GTK_WINDOW(window), config->resizable);
    gtk_window_set_decorated(GTK_WINDOW(window), config->decorated);
    gtk_window_set_modal(GTK_WINDOW(window), config->modal);

    if (config->min_width > 0 || config->min_height > 0) {
        gtk_widget_set_size_request(window, config->min_width, config->min_height);
    }

    apply_window_background(config, window);
    apply_widget_style(window, &config->style);

    if (config->maximized) {
        gtk_window_maximize(GTK_WINDOW(window));
    }
    if (config->present_on_create) {
        gtk_window_present(GTK_WINDOW(window));
    }

    return window;
}
\end{lstlisting}

\subsection{APIs GTK utilisées}

\begin{center}
\begin{tabular}{|l|p{8cm}|}
\hline
\textbf{API} & \textbf{Rôle} \\
\hline
\texttt{gtk\_application\_window\_new()} & Crée une fenêtre associée à une instance de \texttt{GtkApplication}. \\
\hline
\texttt{gtk\_widget\_set\_name()} & Attribue un identifiant CSS au widget pour cibler une seule fenêtre. \\
\hline
\texttt{gtk\_css\_provider\_new()} & Prépare un fournisseur de style CSS chargé dynamiquement. \\
\hline
\texttt{gtk\_style\_context\_add\_provider\_for\_display()} & Enregistre le style CSS au niveau de l'affichage. \\
\hline
\texttt{gtk\_window\_set\_default\_size()} & Définit la taille initiale de la fenêtre. \\
\hline
\texttt{gtk\_window\_present()} & Affiche la fenêtre et demande le focus. \\
\hline
\end{tabular}
\end{center}

\newpage

\section{Module \texttt{button.c}}

Les boutons sont les déclencheurs primaires d'actions. Ce module centralise la logique de création des boutons, qu'ils soient textuels, iconiques ou un mélange des deux. L'apport majeur de ce wrapper est d'intégrer directement le raccordement des signaux (callbacks) lors de la création, éliminant ainsi le risque d'oublier l'appel à \texttt{g\_signal\_connect} plus loin dans le code.

\subsection{Fonction \texttt{create\_button}}

La fonction \texttt{create\_button} est conçue comme un point de rassemblement central pour toutes les formes de déclencheurs interactifs dans GTK. Historiquement, un développeur devait maîtriser différentes fonctions (\texttt{gtk\_button\_new\_with\_label}, \texttt{gtk\_button\_new\_with\_mnemonic}) et jongler avec des conteneurs supplémentaires s'il souhaitait ajouter une icône. Notre approche rationalise drastiquement ce processus.

Dans le corps de la fonction, la première étape logique consiste à discriminer le type de texte demandé. L'utilisation du drapeau booléen \texttt{use\_underline} oriente dynamiquement l'appel vers la fonction de création appropriée, permettant ainsi de supporter nativement les raccourcis clavier (mnémoniques) cruciaux pour l'accessibilité. 

La véritable complexité encapsulée par cette fonction apparaît lors de la gestion des icônes conjointes au texte. Si le champ \texttt{icon\_name} est renseigné, le code bascule d'une simple création textuelle vers la construction transparente d'une hiérarchie de widgets. Il instancie silencieusement une \texttt{GtkBox} interne, y insère une \texttt{GtkImage} configurée avec la taille précise exigée par l'API système, ajoute un \texttt{GtkLabel}, et intègre le tout comme unique enfant du bouton via \texttt{gtk\_button\_set\_child}. 

Cette automatisation cache complètement les rouages de la composition d'interface à l'utilisateur final. Enfin, la fonction s'assure d'injecter immédiatement les données métier (\texttt{user\_data}) et la fonction de rappel (\texttt{on\_clicked}) au signal d'activation, bouclant ainsi la création visuelle et fonctionnelle en un seul appel unifié.

\begin{lstlisting}[language=C]
GtkWidget *create_button(const button_config *config) {
    GtkWidget *button;
    const char *label;

    if (config == NULL) {
        return gtk_button_new();
    }

    label = config->label ? config->label : "";
    button = config->use_underline ? gtk_button_new_with_mnemonic(label) : gtk_button_new_with_label(label);
    gtk_button_set_has_frame(GTK_BUTTON(button), config->has_frame);
    gtk_button_set_can_shrink(GTK_BUTTON(button), config->can_shrink);

    if (config->icon_name && config->icon_name[0] != '\0') {
        GtkWidget *content = gtk_box_new(GTK_ORIENTATION_HORIZONTAL, 6);
        GtkWidget *icon = gtk_image_new_from_icon_name(config->icon_name);
        GtkWidget *text = config->use_underline ? gtk_label_new_with_mnemonic(label) : gtk_label_new(label);
        if (config->icon_size > 0) {
            gtk_image_set_pixel_size(GTK_IMAGE(icon), config->icon_size);
        }
        gtk_box_append(GTK_BOX(content), icon);
        gtk_box_append(GTK_BOX(content), text);
        gtk_button_set_child(GTK_BUTTON(button), content);
    }

    apply_theme_variant(button, config->theme_variant);
    apply_widget_style(button, &config->style);

    if (config->on_clicked) {
        g_signal_connect(button, "clicked", G_CALLBACK(config->on_clicked), config->user_data);
    }

    return button;
}
\end{lstlisting}

\paragraph{Macros associées :}
Les macros de variantes de thèmes (\texttt{BTN\_PRIMARY}, etc.) permettent de s'abstraire des chaînes de caractères brutes lors de la configuration. Bien que simples, elles jouent un rôle crucial pour la fiabilité du code : le compilateur signalera immédiatement une faute de frappe sur le nom de la macro, là où une erreur dans une chaîne littérale comme \texttt{"primay"} passerait inaperçue jusqu'à l'exécution, provoquant un rendu silencieusement incorrect de l'interface. Elles garantissent ainsi l'utilisation de noms de classes CSS cohérents avec la logique interne du module.

\begin{lstlisting}[language=C]
#define BTN_PRIMARY "primary"
#define BTN_SUCCESS "success"
#define BTN_DANGER  "danger"
#define BTN_WARNING "warning"
\end{lstlisting}

\subsection{Fonction \texttt{create\_radio\_button}}

Historiquement complexes à gérer, les boutons radio sont ici simplifiés par une interface déclarative directe. L'implémentation de \texttt{create\_radio\_button} reflète le nouveau standard de GTK4 qui a fusionné les anciens \texttt{GtkRadioButton} au profit d'un composant universel \texttt{GtkCheckButton} gérant des groupes logiques.

Le cœur de cette fonction repose sur un mécanisme d'enchaînement élégant. Au lieu de forcer le développeur à gérer lui-même une liste chaînée externe ou un objet groupe abstrait (comme le \texttt{GSList} sous GTK3), la configuration accepte simplement un pointeur vers un bouton précédemment créé (via le champ \texttt{group\_with}). La fonction vérifie la validité de ce pointeur et, si confirmé, utilise \texttt{gtk\_check\_button\_set\_group} pour raccorder mathématiquement le nouveau widget au groupe d'exclusion existant.

Ce paradigme réduit drastiquement les erreurs logiques de l'utilisateur. De plus, la définition immédiate de l'état par défaut (le booléen \texttt{is\_active}) et le branchement instantané de la callback sur l'événement \texttt{toggled} complètent l'initialisation. Le développeur n'a plus qu'à déclarer ses options l'une après l'autre, laissant le wrapper tisser les liens d'exclusivité en arrière-plan.

\begin{lstlisting}[language=C]
GtkWidget *create_radio_button(const radio_button_config *config) {
    GtkWidget *button = gtk_check_button_new_with_label(config && config->label ? config->label : "");

    if (config == NULL) {
        return button;
    }

    if (config->group_with != NULL) {
        gtk_check_button_set_group(GTK_CHECK_BUTTON(button), GTK_CHECK_BUTTON(config->group_with));
    }
    gtk_check_button_set_active(GTK_CHECK_BUTTON(button), config->is_active);
    apply_widget_style(button, &config->style);

    if (config->on_toggled) {
        g_signal_connect(button, "toggled", G_CALLBACK(config->on_toggled), config->user_data);
    }

    return button;
}
\end{lstlisting}

\subsection{APIs GTK utilisées}

\begin{itemize}
    \item \textbf{\texttt{gtk\_button\_new\_with\_label()}} : Crée un nouveau widget bouton contenant un label texte. Dans le module, elle est utilisée pour instancier le bouton de base lorsque l'option \texttt{use\_underline} est désactivée, offrant une création simplifiée de boutons textuels standards.
    \item \textbf{\texttt{gtk\_button\_new\_with\_mnemonic()}} : Crée un bouton dont le label contient un mnémonique (raccourci clavier) indiqué par un caractère souligné. Le wrapper y a recours lorsque \texttt{config->use\_underline} est vrai, facilitant ainsi l'accessibilité de l'interface au clavier.
    \item \textbf{\texttt{gtk\_button\_set\_child()}} : Définit le widget enfant unique du bouton. Elle permet au wrapper d'injecter une \texttt{GtkBox} personnalisée contenant à la fois une icône et un texte, dépassant ainsi la limitation des boutons à simple label.
    \item \textbf{\texttt{gtk\_image\_new\_from\_icon\_name()}} : Génère un widget d'image à partir du nom d'une icône présente dans le thème système. Elle est utilisée ici pour charger visuellement l'icône demandée dans la configuration avant son insertion dans le bouton.
    \item \textbf{\texttt{gtk\_check\_button\_set\_group()}} : Lie un bouton de type case à cocher à un groupe existant pour obtenir un comportement de sélection exclusive. Dans \texttt{create\_radio\_button}, cette API transforme un \texttt{GtkCheckButton} standard en un bouton radio fonctionnel.
    \item \textbf{\texttt{g\_signal\_connect()}} : Connecte une fonction de rappel à un signal spécifique d'un objet. Le module l'utilise pour brancher les signaux \texttt{clicked} (pour les boutons) et \texttt{toggled} (pour les radios) aux fonctions définies par le développeur applicatif.
\end{itemize}

\newpage

\section{Module \texttt{common.c}}

Ce module constitue l'épine dorsale visuelle du projet. Au lieu de laisser le développeur réécrire les mêmes blocs de code pour chaque widget (marges, expansion, alignement), \texttt{common.c} factorise toutes les propriétés héritées de la classe \texttt{GtkWidget} de GTK. Cela permet de garantir que l'application d'un style est unifiée et prédictible, quel que soit le composant instancié.

\subsection{Fonction \texttt{apply\_widget\_style}}

Cette fonction agit comme un intercepteur universel. Au lieu d'appeler manuellement \texttt{gtk\_widget\_set\_margin\_*} pour chaque bouton ou label, le développeur passe une sous-structure \texttt{style}. L'implémentation est conçue de manière défensive : elle vérifie méthodiquement la présence de pointeurs non nuls et évalue des drapeaux booléens (comme \texttt{set\_halign}) pour garantir qu'aucune propriété GTK n'est écrasée involontairement si elle n'a pas été explicitement demandée par l'utilisateur.

\begin{lstlisting}[language=C]
void apply_widget_style(GtkWidget *widget, const widget_style_config *style) {
    if (widget == NULL || style == NULL) {
        return;
    }

    if (style->css_class && style->css_class[0] != '\0') {
        gtk_widget_add_css_class(widget, style->css_class);
    }
    if (style->css_classes != NULL) {
        for (int i = 0; style->css_classes[i] != NULL; i++) {
            gtk_widget_add_css_class(widget, style->css_classes[i]);
        }
    }

    if (style->width_request != 0 || style->height_request != 0) {
        gtk_widget_set_size_request(widget, style->width_request, style->height_request);
    }

    gtk_widget_set_margin_top(widget, style->margin_top);
    gtk_widget_set_margin_bottom(widget, style->margin_bottom);
    gtk_widget_set_margin_start(widget, style->margin_start);
    gtk_widget_set_margin_end(widget, style->margin_end);

    if (style->set_halign) {
        gtk_widget_set_halign(widget, style->halign);
    }
    if (style->set_valign) {
        gtk_widget_set_valign(widget, style->valign);
    }
    if (style->set_hexpand) {
        gtk_widget_set_hexpand(widget, style->hexpand);
    }
    if (style->set_vexpand) {
        gtk_widget_set_vexpand(widget, style->vexpand);
    }
    if (style->set_sensitive) {
        gtk_widget_set_sensitive(widget, style->sensitive);
    }
    if (style->set_visible) {
        gtk_widget_set_visible(widget, style->visible);
    }

    if (style->tooltip && style->tooltip[0] != '\0') {
        gtk_widget_set_tooltip_text(widget, style->tooltip);
    }
}
\end{lstlisting}

\paragraph{Macros associées :}
L'implémentation de la configuration visuelle est complétée par des macros de substitution. La première série mappe nos propres constantes (ex: \texttt{ALIGN\_CENTER}) directement sur les types énumérés internes de GTK. Cette approche encapsule le code utilisateur en le découplant légèrement de l'API sous-jacente tout en conservant une résolution à la compilation, assurant une performance maximale sans coût d'exécution. La seconde série (\texttt{MARGIN\_*} ) instaure un système de design (Design System) par paliers fixes. En forçant l'utilisation de ces constantes plutôt que des valeurs magiques éparpillées (magic numbers), le wrapper garantit mécaniquement une harmonie spatiale et une cohérence visuelle dans l'ensemble de l'application.

\begin{lstlisting}[language=C]
#define ALIGN_START  GTK_ALIGN_START
#define ALIGN_CENTER GTK_ALIGN_CENTER
#define ALIGN_END    GTK_ALIGN_END
#define ALIGN_FILL   GTK_ALIGN_FILL

#define MARGIN_TINY   4
#define MARGIN_SMALL  8
#define MARGIN_MEDIUM 16
#define MARGIN_LARGE  24
\end{lstlisting}

\subsection{APIs GTK utilisées}

\begin{itemize}
    \item \textbf{\texttt{gtk\_widget\_add\_css\_class()}} : Ajoute une classe de style CSS spécifiée au widget cible. Dans \texttt{apply\_widget\_style}, elle permet d'appliquer de manière dynamique une ou plusieurs classes CSS définies dans la configuration pour modifier l'apparence visuelle du widget.
    \item \textbf{\texttt{gtk\_widget\_set\_size\_request()}} : Fixe les dimensions minimales (largeur et hauteur) que le widget doit tenter d'occuper. Le wrapper utilise cette fonction pour forcer une taille spécifique demandée par le développeur dans la structure \texttt{widget\_style\_config}.
    \item \textbf{\texttt{gtk\_widget\_set\_margin\_top()}} (et variantes) : Définit les marges extérieures du widget pour chaque direction. Le module les exploite pour assurer un espacement cohérent entre les widgets sans avoir à manipuler directement les conteneurs parents.
    \item \textbf{\texttt{gtk\_widget\_set\_halign()}} et \textbf{\texttt{gtk\_widget\_set\_valign()}} : Contrôlent l'alignement du widget respectivement sur les axes horizontal et vertical. Elles permettent au wrapper de gérer la position du widget (centré, aligné à gauche/droite, ou étendu) dans sa zone d'affichage.
    \item \textbf{\texttt{gtk\_widget\_set\_hexpand()}} et \textbf{\texttt{gtk\_widget\_set\_vexpand()}} : Indiquent au gestionnaire de mise en page si le widget doit absorber l'espace supplémentaire disponible. Ces appels sont cruciaux pour définir le comportement de redimensionnement réactif des interfaces construites avec le wrapper.
    \item \textbf{\texttt{gtk\_widget\_set\_tooltip\_text()}} : Associe un texte d'aide contextuelle qui apparaît lorsque l'utilisateur survole le widget avec la souris. Elle est systématiquement appelée par le module si une chaîne de caractères \texttt{tooltip} est fournie.
\end{itemize}

\newpage

\section{Module \texttt{container.c}}

Le module \texttt{container.c} rassemble les conteneurs de mise en page du wrapper. Il couvre la grille, la boîte linéaire, la pile de pages et deux helpers complémentaires pour les interfaces multi-vues.

\subsection{Fonction \texttt{create\_grid}}

La fonction \texttt{create\_grid} instancie un \texttt{GtkGrid} puis applique l'espacement entre lignes et colonnes ainsi que les options d'homogénéité.

\begin{lstlisting}[language=C]
GtkWidget *create_grid(const grid_config *config) {
    GtkWidget *grid = gtk_grid_new();
    if (config == NULL) {
        return grid;
    }

    gtk_grid_set_column_spacing(GTK_GRID(grid), config->column_spacing);
    gtk_grid_set_row_spacing(GTK_GRID(grid), config->row_spacing);
    gtk_grid_set_column_homogeneous(GTK_GRID(grid), config->homogeneous_columns);
    gtk_grid_set_row_homogeneous(GTK_GRID(grid), config->homogeneous_rows);
    apply_widget_style(grid, &config->style);
    return grid;
}
\end{lstlisting}

\subsection{Fonction \texttt{create\_box}}

\begin{lstlisting}[language=C]
GtkWidget *create_box(const box_config *config) {
    GtkWidget *box = gtk_box_new(
        config ? config->orientation : GTK_ORIENTATION_VERTICAL,
        config ? config->spacing : 0
    );
    if (config == NULL) {
        return box;
    }

    gtk_box_set_homogeneous(GTK_BOX(box), config->homogeneous);
    apply_widget_style(box, &config->style);
    return box;
}
\end{lstlisting}

\subsection{Fonction \texttt{create\_stack}}

\begin{lstlisting}[language=C]
GtkWidget *create_stack(const stack_config *config) {
    GtkWidget *stack = gtk_stack_new();
    if (config == NULL) {
        return stack;
    }

    gtk_stack_set_transition_type(GTK_STACK(stack), config->transition);
    gtk_stack_set_transition_duration(GTK_STACK(stack), config->duration_ms);
    gtk_stack_set_vhomogeneous(GTK_STACK(stack), config->vhomogeneous);
    gtk_stack_set_hhomogeneous(GTK_STACK(stack), config->hhomogeneous);
    apply_widget_style(stack, &config->style);
    return stack;
}
\end{lstlisting}

\subsection{Helpers complémentaires}

Le module ajoute aussi deux helpers pratiques : \texttt{create\_stack\_sidebar}, qui relie une barre latérale à une pile existante, et \texttt{stack\_add\_page}, qui ajoute une page puis renseigne son titre par propriété GTK.

\begin{lstlisting}[language=C]
GtkWidget *create_stack_sidebar(GtkWidget *stack, const widget_style_config *style) {
    GtkWidget *sidebar = gtk_stack_sidebar_new();
    if (stack != NULL) {
        gtk_stack_sidebar_set_stack(GTK_STACK_SIDEBAR(sidebar), GTK_STACK(stack));
    }
    apply_widget_style(sidebar, style);
    return sidebar;
}

void stack_add_page(GtkWidget *stack, GtkWidget *child, const char *name, const char *title) {
    if (stack == NULL || child == NULL || name == NULL || title == NULL) {
        return;
    }

    gtk_stack_add_named(GTK_STACK(stack), child, name);
    g_object_set(gtk_stack_get_page(GTK_STACK(stack), child), "title", title, NULL);
}
\end{lstlisting}

\subsection{APIs GTK utilisées}

\begin{center}
\begin{tabular}{|l|p{8cm}|}
\hline
\textbf{API} & \textbf{Rôle} \\
\hline
\texttt{gtk\_grid\_new()} & Crée un conteneur en grille. \\
\hline
\texttt{gtk\_box\_new()} & Crée un conteneur linéaire horizontal ou vertical. \\
\hline
\texttt{gtk\_stack\_new()} & Crée un conteneur à pages multiples. \\
\hline
\texttt{gtk\_stack\_sidebar\_new()} & Crée le widget de navigation latérale lié à une pile. \\
\hline
\texttt{gtk\_stack\_add\_named()} & Ajoute une page identifiée dans la pile. \\
\hline
\texttt{g\_object\_set()} & Attribue le titre de page utilisé par les widgets de navigation. \\
\hline
\end{tabular}
\end{center}

\newpage

\section{Module \texttt{display.c}}

Le module \texttt{display.c} prend en charge les widgets d'affichage pur : barre de progression, spinner, label et image. Son rôle est d'exposer des constructeurs simples au-dessus des widgets GTK standards.

\subsection{Fonction \texttt{create\_progress\_bar}}

\begin{lstlisting}[language=C]
GtkWidget *create_progress_bar(const progress_bar_config *config) {
    GtkWidget *progress = gtk_progress_bar_new();

    if (config == NULL) {
        return progress;
    }

    gtk_progress_bar_set_fraction(GTK_PROGRESS_BAR(progress), config->fraction);
    gtk_progress_bar_set_show_text(GTK_PROGRESS_BAR(progress), config->show_text);
    gtk_progress_bar_set_inverted(GTK_PROGRESS_BAR(progress), config->inverted);
    gtk_orientable_set_orientation(GTK_ORIENTABLE(progress), config->orientation);
    if (config->text != NULL) {
        gtk_progress_bar_set_text(GTK_PROGRESS_BAR(progress), config->text);
    }
    if (config->pulse_once) {
        gtk_progress_bar_pulse(GTK_PROGRESS_BAR(progress));
    }
    apply_widget_style(progress, &config->style);

    return progress;
}
\end{lstlisting}

\subsection{Fonction \texttt{create\_spinner}}

\begin{lstlisting}[language=C]
GtkWidget *create_spinner(const spinner_config *config) {
    GtkWidget *spinner = gtk_spinner_new();

    if (config != NULL) {
        if (config->size > 0) {
            gtk_widget_set_size_request(spinner, config->size, config->size);
        }
        if (config->spinning) {
            gtk_spinner_start(GTK_SPINNER(spinner));
        }
        apply_widget_style(spinner, &config->style);
    }

    return spinner;
}
\end{lstlisting}

\subsection{Fonction \texttt{create\_label}}

\begin{lstlisting}[language=C]
GtkWidget *create_label(const label_config *config) {
    GtkWidget *label = gtk_label_new(config && config->text ? config->text : "");

    if (config == NULL) {
        return label;
    }

    gtk_label_set_selectable(GTK_LABEL(label), config->selectable);
    gtk_label_set_wrap(GTK_LABEL(label), config->wrap);
    gtk_label_set_ellipsize(GTK_LABEL(label), config->ellipsize);
    if (config->max_width_chars > 0) {
        gtk_label_set_max_width_chars(GTK_LABEL(label), config->max_width_chars);
    }
    apply_widget_style(label, &config->style);

    return label;
}
\end{lstlisting}

\subsection{Fonction \texttt{create\_image}}

\begin{lstlisting}[language=C]
GtkWidget *create_image(const image_config *config) {
    GtkWidget *widget;

    if (config == NULL) {
        return gtk_image_new();
    }

    if (config->file_path != NULL && config->file_path[0] != '\0') {
        GtkWidget *pic = gtk_picture_new_for_filename(config->file_path);
        gtk_picture_set_can_shrink(GTK_PICTURE(pic), TRUE);
        gtk_picture_set_content_fit(
            GTK_PICTURE(pic),
            config->keep_aspect ? GTK_CONTENT_FIT_CONTAIN : GTK_CONTENT_FIT_FILL
        );
        widget = pic;
    } else if (config->icon_name != NULL && config->icon_name[0] != '\0') {
        GtkWidget *img = gtk_image_new_from_icon_name(config->icon_name);
        if (config->pixel_size > 0) {
            gtk_image_set_pixel_size(GTK_IMAGE(img), config->pixel_size);
        }
        widget = img;
    } else {
        widget = gtk_image_new();
    }

    apply_widget_style(widget, &config->style);
    return widget;
}
\end{lstlisting}

\subsection{APIs GTK utilisées}

\begin{center}
\begin{tabular}{|l|p{8cm}|}
\hline
\textbf{API} & \textbf{Rôle} \\
\hline
\texttt{gtk\_progress\_bar\_set\_fraction()} & Positionne l'avancement sur une valeur normalisée. \\
\hline
\texttt{gtk\_progress\_bar\_pulse()} & Déclenche un avancement indéterminé ponctuel. \\
\hline
\texttt{gtk\_spinner\_start()} & Lance l'animation du spinner. \\
\hline
\texttt{gtk\_label\_set\_ellipsize()} & Contrôle la coupure des textes trop longs. \\
\hline
\texttt{gtk\_picture\_new\_for\_filename()} & Crée un widget d'image basé sur un fichier. \\
\hline
\texttt{gtk\_image\_new\_from\_icon\_name()} & Crée une icône issue du thème GTK. \\
\hline
\end{tabular}
\end{center}

\newpage

\section{Module \texttt{input.c}}

Le module \texttt{input.c} regroupe les composants de saisie. Il encapsule trois widgets principaux : l'entrée de texte, la liste déroulante moderne et le bouton numérique à incréments.

\subsection{Fonction \texttt{create\_entry}}

Cette fonction prépare un \texttt{GtkEntry} puis applique le placeholder, le texte initial, le type de saisie attendu et les éventuelles contraintes de visibilité ou de longueur.

\begin{lstlisting}[language=C]
GtkWidget *create_entry(const entry_config *config) {
    GtkWidget *entry = gtk_entry_new();

    if (config == NULL) {
        return entry;
    }

    if (config->placeholder) {
        gtk_entry_set_placeholder_text(GTK_ENTRY(entry), config->placeholder);
    }
    if (config->default_text) {
        gtk_editable_set_text(GTK_EDITABLE(entry), config->default_text);
    }
    gtk_entry_set_input_purpose(GTK_ENTRY(entry), parse_input_purpose(config->purpose_hint));
    gtk_entry_set_visibility(GTK_ENTRY(entry), !config->password);
    gtk_editable_set_editable(GTK_EDITABLE(entry), config->editable);
    if (config->max_length > 0) {
        gtk_entry_set_max_length(GTK_ENTRY(entry), config->max_length);
    }
    apply_widget_style(entry, &config->style);

    if (config->on_changed) {
        g_signal_connect(entry, "changed", G_CALLBACK(config->on_changed), config->user_data);
    }

    return entry;
}
\end{lstlisting}

\subsection{Fonction \texttt{create\_dropdown}}

\begin{lstlisting}[language=C]
GtkWidget *create_dropdown(const dropdown_config *config) {
    GtkWidget *dropdown;

    if (config == NULL || config->options == NULL) {
        return gtk_drop_down_new(NULL, NULL);
    }

    dropdown = gtk_drop_down_new_from_strings(config->options);
    gtk_drop_down_set_selected(GTK_DROP_DOWN(dropdown), config->selected_index);
    gtk_drop_down_set_enable_search(GTK_DROP_DOWN(dropdown), config->enable_search);
    apply_widget_style(dropdown, &config->style);

    if (config->on_selected) {
        g_signal_connect(dropdown, "notify::selected", G_CALLBACK(config->on_selected), config->user_data);
    }

    return dropdown;
}
\end{lstlisting}

\subsection{Fonction \texttt{create\_spin\_button}}

\begin{lstlisting}[language=C]
GtkWidget *create_spin_button(const spin_button_config *config) {
    GtkAdjustment *adj;
    GtkWidget *spin;

    if (config == NULL) {
        adj = gtk_adjustment_new(0.0, 0.0, 100.0, 1.0, 5.0, 0.0);
        return gtk_spin_button_new(adj, 1.0, 0);
    }

    adj = gtk_adjustment_new(config->value, config->min_value, config->max_value, config->step, config->page, 0.0);
    spin = gtk_spin_button_new(adj, config->step, config->digits);
    gtk_spin_button_set_numeric(GTK_SPIN_BUTTON(spin), config->numeric_only);
    gtk_spin_button_set_wrap(GTK_SPIN_BUTTON(spin), config->wrap);
    apply_widget_style(spin, &config->style);

    if (config->on_value_changed) {
        g_signal_connect(spin, "value-changed", G_CALLBACK(config->on_value_changed), config->user_data);
    }

    return spin;
}
\end{lstlisting}

\subsection{APIs GTK utilisées}

\begin{center}
\begin{tabular}{|l|p{8cm}|}
\hline
\textbf{API} & \textbf{Rôle} \\
\hline
\texttt{gtk\_entry\_set\_placeholder\_text()} & Définit le texte d'indication d'un champ vide. \\
\hline
\texttt{gtk\_entry\_set\_input\_purpose()} & Informe GTK du type de saisie attendu. \\
\hline
\texttt{gtk\_drop\_down\_new\_from\_strings()} & Construit une liste de choix depuis un tableau de chaînes. \\
\hline
\texttt{gtk\_drop\_down\_set\_enable\_search()} & Active la recherche intégrée du widget déroulant. \\
\hline
\texttt{gtk\_adjustment\_new()} & Définit l'intervalle numérique et les pas du contrôle. \\
\hline
\texttt{gtk\_spin\_button\_new()} & Crée le widget de saisie numérique associé à un ajustement. \\
\hline
\end{tabular}
\end{center}

\newpage

\section{Module \texttt{toggle.c}}

Le module \texttt{toggle.c} couvre les composants binaires du wrapper : case à cocher et ligne composée avec interrupteur. Ces deux widgets reposent sur des contrôles GTK simples, mais le wrapper leur ajoute une configuration déclarative homogène.

\subsection{Fonction \texttt{create\_checkbox}}

\begin{lstlisting}[language=C]
GtkWidget *create_checkbox(const checkbox_config *config) {
    GtkWidget *checkbox = gtk_check_button_new_with_label(config && config->label ? config->label : "");

    if (config == NULL) {
        return checkbox;
    }

    gtk_check_button_set_active(GTK_CHECK_BUTTON(checkbox), config->active);
    gtk_check_button_set_inconsistent(GTK_CHECK_BUTTON(checkbox), config->inconsistent);
    apply_widget_style(checkbox, &config->style);

    if (config->on_toggled != NULL) {
        g_signal_connect(checkbox, "toggled", G_CALLBACK(config->on_toggled), config->user_data);
    }

    return checkbox;
}
\end{lstlisting}

\subsection{Fonction \texttt{create\_switch\_row}}

Cette seconde fonction construit un composant composite : un label à gauche, un \texttt{GtkSwitch} à droite, puis un éventuel callback sur la propriété \texttt{active}.

\begin{lstlisting}[language=C]
GtkWidget *create_switch_row(const switch_config *config, GtkWidget **out_switch) {
    GtkWidget *row = gtk_box_new(GTK_ORIENTATION_HORIZONTAL, 10);
    GtkWidget *label = gtk_label_new(config && config->label ? config->label : "");
    GtkWidget *sw = gtk_switch_new();

    gtk_widget_set_hexpand(label, TRUE);
    gtk_widget_set_halign(label, GTK_ALIGN_START);
    gtk_widget_set_halign(sw, GTK_ALIGN_END);
    gtk_box_append(GTK_BOX(row), label);
    gtk_box_append(GTK_BOX(row), sw);

    if (config != NULL) {
        gtk_switch_set_active(GTK_SWITCH(sw), config->active);
        gtk_switch_set_state(GTK_SWITCH(sw), config->state);
        apply_widget_style(row, &config->style);

        if (config->on_active_changed != NULL) {
            g_signal_connect(sw, "notify::active", G_CALLBACK(config->on_active_changed), config->user_data);
        }
    }

    if (out_switch != NULL) {
        *out_switch = sw;
    }

    return row;
}
\end{lstlisting}

\subsection{APIs GTK utilisées}

\begin{center}
\begin{tabular}{|l|p{8cm}|}
\hline
\textbf{API} & \textbf{Rôle} \\
\hline
\texttt{gtk\_check\_button\_new\_with\_label()} & Crée la case à cocher avec son libellé. \\
\hline
\texttt{gtk\_check\_button\_set\_inconsistent()} & Active l'état visuel intermédiaire. \\
\hline
\texttt{gtk\_switch\_new()} & Crée l'interrupteur binaire GTK4. \\
\hline
\texttt{gtk\_switch\_set\_active()} & Positionne l'état logique initial du switch. \\
\hline
\texttt{gtk\_box\_append()} & Assemble le label et l'interrupteur dans une même ligne. \\
\hline
\texttt{g\_signal\_connect()} & Relie les changements d'état aux callbacks utilisateur. \\
\hline
\end{tabular}
\end{center}

\newpage

\section{Module \texttt{date.c}}

Le module \texttt{date.c} encapsule le widget calendrier de GTK4. Il prend en charge la date initiale, l'affichage des noms de jours, les numéros de semaine et la remontée du signal de sélection.

\subsection{Fonction \texttt{create\_calendar}}

\begin{lstlisting}[language=C]
GtkWidget *create_calendar(const calendar_config *config) {
    GtkWidget *calendar = gtk_calendar_new();

    if (config == NULL) {
        return calendar;
    }

    if (config->selected_date != NULL) {
#if GTK_CHECK_VERSION(4, 20, 0)
        gtk_calendar_set_date(GTK_CALENDAR(calendar), config->selected_date);
#else
        gtk_calendar_select_day(GTK_CALENDAR(calendar), config->selected_date);
#endif
    }
    gtk_calendar_set_show_day_names(GTK_CALENDAR(calendar), config->show_day_names);
    gtk_calendar_set_show_week_numbers(GTK_CALENDAR(calendar), config->show_week_numbers);
    gtk_calendar_set_show_heading(GTK_CALENDAR(calendar), !config->no_month_change);
    apply_widget_style(calendar, &config->style);

    if (config->on_day_selected != NULL) {
        g_signal_connect(calendar, "day-selected", G_CALLBACK(config->on_day_selected), config->user_data);
    }

    return calendar;
}
\end{lstlisting}

\subsection{APIs GTK utilisées}

\begin{center}
\begin{tabular}{|l|p{8cm}|}
\hline
\textbf{API} & \textbf{Rôle} \\
\hline
\texttt{gtk\_calendar\_new()} & Crée l'instance du calendrier. \\
\hline
\texttt{gtk\_calendar\_set\_date()} & Positionne la date courante dans les versions récentes de GTK. \\
\hline
\texttt{gtk\_calendar\_select\_day()} & Assure la compatibilité avec les versions plus anciennes. \\
\hline
\texttt{gtk\_calendar\_set\_show\_week\_numbers()} & Active l'affichage du numéro de semaine. \\
\hline
\texttt{gtk\_calendar\_set\_show\_heading()} & Contrôle l'en-tête et la navigation mensuelle. \\
\hline
\texttt{g\_signal\_connect()} & Raccorde le signal de sélection de jour au callback. \\
\hline
\end{tabular}
\end{center}

\newpage

\section{Module \texttt{menu.c}}

Le module \texttt{menu.c} implémente l'architecture moderne des menus GTK4. Les éléments visibles sont décrits via \texttt{GMenuModel}, tandis que la logique applicative est fournie sous forme d'actions enregistrées dans un \texttt{GActionMap}.

\subsection{Fonction \texttt{register\_item\_action}}

Cette fonction privée crée une action GTK pour une entrée de menu active, connecte le callback utilisateur puis enregistre l'action dans la map de l'application.

\begin{lstlisting}[language=C]
static void register_item_action(const menubar_config *config,
                                 const menu_item_config *item,
                                 const char *action_name) {
    if (config == NULL || config->action_map == NULL || item == NULL || item->on_activate == NULL || action_name == NULL) {
        return;
    }

    GSimpleAction *action = g_simple_action_new(action_name, NULL);
    g_signal_connect(action, "activate", G_CALLBACK(item->on_activate), item->user_data);
    g_action_map_add_action(config->action_map, G_ACTION(action));
    g_object_unref(action);
}
\end{lstlisting}

\subsection{Fonction \texttt{append\_item\_recursive}}

La construction du menu est récursive. Le code traite d'abord les sous-menus et les sections, puis les actions simples en leur attribuant un nom unique.

\begin{lstlisting}[language=C]
static void append_item_recursive(GMenu *menu,
                                  const menu_item_config *item,
                                  const menubar_config *config) {
    if (item == NULL || item->enabled == false) {
        return;
    }

    if (item->kind == MENU_ITEM_SUBMENU && item->children != NULL && item->child_count > 0) {
        GMenu *submenu = g_menu_new();
        for (size_t i = 0; i < item->child_count; i++) {
            append_item_recursive(submenu, &item->children[i], config);
        }
        g_menu_append_submenu(menu, item->label ? item->label : "", G_MENU_MODEL(submenu));
        g_object_unref(submenu);
        return;
    }

    if (item->kind == MENU_ITEM_SECTION && item->children != NULL && item->child_count > 0) {
        GMenu *section = g_menu_new();
        for (size_t i = 0; i < item->child_count; i++) {
            append_item_recursive(section, &item->children[i], config);
        }
        g_menu_append_section(menu, item->label, G_MENU_MODEL(section));
        g_object_unref(section);
        return;
    }

    char *action_name = create_action_name();
    register_item_action(config, item, action_name);

    char *detailed_action = g_strdup_printf("app.%s", action_name);
    g_menu_append(menu, item->label ? item->label : "", detailed_action);

    g_free(detailed_action);
    g_free(action_name);
}
\end{lstlisting}

\subsection{Fonctions \texttt{create\_menubar} et \texttt{create\_menu\_panel}}

Les fonctions publiques consomment un modèle construit par \texttt{build\_menu\_model}. La première fabrique une barre de menus, la seconde un panneau de menu réutilisable dans un popover.

\begin{lstlisting}[language=C]
GtkWidget *create_menubar(const menubar_config *config) {
    GMenuModel *model = build_menu_model(config);
    GtkWidget *bar = gtk_popover_menu_bar_new_from_model(model);
    if (config != NULL) {
        apply_widget_style(bar, &config->style);
    }
    g_object_unref(model);
    return bar;
}

GtkWidget *create_menu_panel(const menubar_config *config) {
    GMenuModel *model = build_menu_model(config);
    GtkWidget *panel = gtk_popover_menu_new_from_model(model);

    if (config != NULL) {
        gtk_orientable_set_orientation(
            GTK_ORIENTABLE(panel),
            config->layout == MENU_LAYOUT_VERTICAL ? GTK_ORIENTATION_VERTICAL : GTK_ORIENTATION_HORIZONTAL
        );
        gtk_popover_set_has_arrow(GTK_POPOVER(panel), config->show_arrow_indicators);
        apply_widget_style(panel, &config->style);
    }

    g_object_unref(model);
    return panel;
}
\end{lstlisting}

\subsection{APIs GTK utilisées}

\begin{center}
\begin{tabular}{|l|p{8cm}|}
\hline
\textbf{API} & \textbf{Rôle} \\
\hline
\texttt{g\_simple\_action\_new()} & Crée l'action logique associée à un item de menu. \\
\hline
\texttt{g\_action\_map\_add\_action()} & Enregistre l'action dans la carte fournie par l'application. \\
\hline
\texttt{g\_menu\_append\_submenu()} & Ajoute un sous-menu hiérarchique au modèle. \\
\hline
\texttt{g\_menu\_append\_section()} & Crée une section logique de menu. \\
\hline
\texttt{gtk\_popover\_menu\_bar\_new\_from\_model()} & Instancie la barre de menus GTK4 à partir du modèle. \\
\hline
\texttt{gtk\_popover\_menu\_new\_from\_model()} & Produit un panneau de menu pouvant être affiché dans un popover. \\
\hline
\end{tabular}
\end{center}

\newpage

\section{Module \texttt{dialog.c}}

Le module \texttt{dialog.c} implémente les boîtes d'alerte asynchrones. Il s'appuie sur \texttt{GtkAlertDialog}, qui ne bloque pas l'exécution du programme et impose donc un traitement différé de la réponse utilisateur.

\subsection{Callback interne \texttt{on\_alert\_response}}

Cette fonction récupère le résultat asynchrone, transforme une éventuelle erreur en indice \texttt{-1}, puis transmet la réponse au callback utilisateur avant de libérer les ressources temporaires.

\begin{lstlisting}[language=C]
static void on_alert_response(GObject *source_object, GAsyncResult *result, gpointer user_data) {
    (void)source_object;
    alert_async_context *ctx = user_data;
    GError *error = NULL;
    int index = gtk_alert_dialog_choose_finish(ctx->dialog, result, &error);
    if (error != NULL) {
        g_error_free(error);
        index = -1;
    }
    if (ctx->on_response != NULL) {
        ctx->on_response(index, ctx->user_data);
    }
    g_object_unref(ctx->dialog);
    g_free(ctx);
}
\end{lstlisting}

\subsection{Fonction \texttt{show\_alert\_dialog}}

La fonction publique configure le message, les détails, la modalité et la liste des boutons. Elle déduit aussi le bouton par défaut et le bouton d'annulation à partir de la configuration avant d'afficher le dialogue en mode synchrone visuel mais asynchrone du point de vue du code.

\begin{lstlisting}[language=C]
void show_alert_dialog(const alert_dialog_config *config) {
    GtkAlertDialog *dialog;
    const char **buttons;

    if (config == NULL || config->message == NULL) {
        return;
    }

    dialog = gtk_alert_dialog_new("%s", config->message);
    if (config->detail != NULL) {
        gtk_alert_dialog_set_detail(dialog, config->detail);
    }
    gtk_alert_dialog_set_modal(dialog, config->modal);

    if (config->button_count > 0 && config->buttons != NULL) {
        buttons = g_new0(const char *, config->button_count + 1);
        for (size_t i = 0; i < config->button_count; i++) {
            buttons[i] = config->buttons[i].label;
        }
        gtk_alert_dialog_set_buttons(dialog, buttons);
        g_free(buttons);
    }

    int default_button = config->default_button;
    int cancel_button = config->cancel_button;

    if (config->buttons != NULL && config->button_count > 0) {
        for (size_t i = 0; i < config->button_count; i++) {
            const char *appearance = config->buttons[i].appearance;
            if (appearance == NULL) {
                continue;
            }
            if (default_button < 0 && strcmp(appearance, "suggested") == 0) {
                default_button = (int)i;
            }
            if (cancel_button < 0 && strcmp(appearance, "cancel") == 0) {
                cancel_button = (int)i;
            }
        }
    }

    if (default_button >= 0) {
        gtk_alert_dialog_set_default_button(dialog, default_button);
    }
    if (cancel_button >= 0) {
        gtk_alert_dialog_set_cancel_button(dialog, cancel_button);
    }

    if (config->on_response != NULL) {
        alert_async_context *ctx = g_new0(alert_async_context, 1);
        ctx->dialog = g_object_ref(dialog);
        ctx->on_response = config->on_response;
        ctx->user_data = config->user_data;
        gtk_alert_dialog_choose(dialog, config->parent, NULL, on_alert_response, ctx);
    } else {
        gtk_alert_dialog_show(dialog, config->parent);
    }

    g_object_unref(dialog);
}
\end{lstlisting}

\subsection{APIs GTK utilisées}

\begin{center}
\begin{tabular}{|l|p{8cm}|}
\hline
\textbf{API} & \textbf{Rôle} \\
\hline
\texttt{gtk\_alert\_dialog\_new()} & Crée le dialogue d'alerte GTK4. \\
\hline
\texttt{gtk\_alert\_dialog\_set\_buttons()} & Définit l'ensemble des boutons proposés à l'utilisateur. \\
\hline
\texttt{gtk\_alert\_dialog\_set\_default\_button()} & Choisit le bouton activé par défaut. \\
\hline
\texttt{gtk\_alert\_dialog\_choose()} & Lance l'affichage asynchrone avec callback de fin. \\
\hline
\texttt{gtk\_alert\_dialog\_choose\_finish()} & Récupère l'indice du bouton choisi dans le résultat asynchrone. \\
\hline
\texttt{gtk\_alert\_dialog\_show()} & Affiche le dialogue sans traitement de réponse personnalisé. \\
\hline
\end{tabular}
\end{center}

\newpage

\section{Module \texttt{separator.c}}

Le module \texttt{separator.c} est volontairement minimal. Il fournit un simple constructeur pour les séparateurs visuels, tout en restant cohérent avec le reste du wrapper grâce à \texttt{widget\_style\_config}.

\subsection{Fonction \texttt{create\_separator}}

\begin{lstlisting}[language=C]
GtkWidget *create_separator(const separator_config *config) {
    GtkWidget *separator = gtk_separator_new(config ? config->orientation : GTK_ORIENTATION_HORIZONTAL);
    if (config != NULL) {
        apply_widget_style(separator, &config->style);
    }
    return separator;
}
\end{lstlisting}

\subsection{APIs GTK utilisées}

\begin{center}
\begin{tabular}{|l|p{8cm}|}
\hline
\textbf{API} & \textbf{Rôle} \\
\hline
\texttt{gtk\_separator\_new()} & Crée un séparateur horizontal ou vertical selon l'orientation demandée. \\
\hline
\end{tabular}
\end{center}

\newpage

\input{demo_app}

\end{document}
